\section{Discussion}
\label{sec:discussion}

\para{Why the main hypothesis fails.} The original prediction was that weaker factual representations in lower-resource languages would make challenge-induced agreement pressure more effective. The factual half of that story is supported: both models are substantially worse on \mkqa in lower-resource languages. The capitulation half is not. On \bullshitbench, lower-resource languages usually move toward stronger rejection, not stronger engagement. The simplest interpretation is that the benchmark is measuring two behaviors at once: factual weakness and conservative false-premise rejection.

\para{A split failure profile.} Our results suggest that multilingual reliability should not be summarized by a single scalar such as ``sycophancy rate.'' Instead, there are at least two relevant failure modes. One is \emph{knowledge weakness}, which clearly appears in the factual QA setting. The other is \emph{challenge response style}, which can become refusal-heavy on nonsense-premise prompts. If a challenge turn doubles as a second-chance reconsideration prompt, then increased rejection can mask any underlying tendency to agree with the user.

\para{What the mechanistic result adds.} The mechanistic follow-up still matters even though the cross-lingual behavioral story is weaker than expected. Qwen contains a compact early-layer direction that separates reject from engage outcomes on challenged \bullshitbench contexts, and moving generation against that direction sharply improves final rejection. This does not prove a language-resource scaling law, but it does show that the challenged response policy is locally steerable and behaviorally consequential.

\para{Limitations.} The study has several clear limitations. First, \bullshitbench is a false-premise benchmark, not a pure pressure-to-agree benchmark, so the challenge turn can increase rejection instead of eliciting capitulation. Second, some multilingual prompts and aliases were translated rather than taken from a fully official benchmark release, which leaves room for translation artifacts. Third, \gptmini serves both as an evaluated model and as the \bullshitbench judge. Fourth, the mechanistic analysis uses a hidden-state contrast heuristic rather than head-level localization or path patching. Fifth, model coverage is narrow: one API model and one open-weight model are enough to test the hypothesis, but not enough to establish broad scaling claims.

\para{Broader implications.} The main practical lesson is not that multilingual sycophancy is absent. It is that English-only safety evaluation can miss a different multilingual risk: lower-resource factual collapse paired with benchmark-specific conservatism. Future work should use challenge prompts that pressure agreement with a specific incorrect answer, compare literal and culturally adapted translations, and test more models so that language effects can be separated from model-family artifacts.
